\section{Geometría, evaluaciones y control del marco}
\label{sec:geometria}

La comparación de la Nota~\ref{sec:periodos} utiliza pendientes medias,
pendientes extremas y una estimación colectiva de los mapas de periodos.
Estas cantidades proceden de construcciones diferentes. Los espacios de
secciones de fibrados proporcionan retículos; las evaluaciones y las
comparaciones de métricas permiten estimar sus pendientes; las correlaciones
entre columnas controlan las dos trazas del operador del marco.
El objetivo de esta nota es demostrar los pasos finitos que convierten
esas entradas en las desigualdades de la Nota~\ref{sec:marcos}.

Se conserva la pendiente $\mu$ de \eqref{eq:marcos-normalizacion},
igual al doble de la pendiente de Arakelov habitual. En A se escribe
$\mu_2$. Los grados son aditivos para una sucesión exacta equipada
con las métricas inducida y cociente. Los índices finitos contribuyen
con el factor $2/[K:\Q]$ de \eqref{eq:marcos-coste-indice}.
Las constantes de esta nota pertenecen a una familia previamente fijada;
su independencia de los índices variables forma parte de cada entrada.

\fuente{A, §7, pp.~28--34, con su enlace a §8; B, §4.3, pp.~33--40,
§4.4, pp.~41--46, y el enlace con §4.5, p.~46. B, §5, identifica el coste
cohomológico que permanece pendiente. Se usan los mismos ejemplares
de 55 y 73 páginas de las notas anteriores.}

\subsection{Secciones de fibrados y retículos de formas}

Un fibrado en rectas $\mathcal L$ asigna a cada punto una recta que
varía algebraica u holomórficamente; localmente sus secciones se escriben
como funciones multiplicadas por un generador. Sobre un modelo propio
$\mathcal X$ definido sobre $\OK$, el módulo
\[
 M_k=H^0(\mathcal X,\mathcal L^{\otimes k})
\]
reúne las secciones integrales. Cuando es finito y libre de torsión,
su fibra sobre $K$ es un espacio vectorial y $M_k$ es un retículo en él.
Una métrica hermitiana en el fibrado produce normas sobre las secciones
en cada inmersión compleja, por ejemplo la norma supremo o una norma
$L^2$ con una medida especificada. Comparar ambas normas exige una
estimación; sus covolúmenes no se identifican sin ese paso.

En B, pp.~33--34, las formas automorfas holomorfas se realizan mediante
fibrados sobre modelos de espacios de módulos de variedades abelianas.
Si $\pi:\mathcal A\to\mathcal M$ es la familia universal y $e$ su
sección cero, el fibrado de Hodge es
$\omega=e^*\Omega^1_{\mathcal A/\mathcal M}$.
La acción del anillo de enteros del cuerpo CM lo descompone, tras
escoger un cuerpo de definición adecuado. Los pesos especifican
potencias de rectas obtenidas de esos sumandos y de sus banderas.
La condición cuspidal se expresa mediante anulación en la frontera
de una compactificación.

Esta descripción identifica los objetos que intervienen en las
estimaciones. Para aplicarla se necesitan modelos integrales,
extensión de los fibrados a la frontera, comparaciones de normas
y control de los cambios de peso y nivel. La realización compleja
de las formas no proporciona por sí sola esos datos finitos.

\subsection{Secciones pequeñas, generación y pendiente}

Para un fibrado hermitiano $\overline{\mathcal L}$ sobre una variedad
aritmética de dimensión total $d$, el conteo de secciones pequeñas es
\[
 \widehat h^0(k\overline{\mathcal L})
 =\log\#\{s\in H^0(\mathcal X,\mathcal L^{\otimes k}):
                       \|s\|_{\sup}\le1\}.
\]
Su crecimiento define el volumen aritmético
\[
 \widehat{\operatorname{vol}}(\overline{\mathcal L})
 =\limsup_{k\to\infty}\frac{d!}{k^d}
                  \widehat h^0(k\overline{\mathcal L}).
\]
B, pp.~34--36, distingue este conteo de la característica
$\chi_{\sup}$ definida mediante volumen y covolumen. Para un fibrado
aritméticamente nef con métricas suaves, el teorema de Hilbert--Samuel de
Moriwaki identifica el término principal del conteo con la
autointersección aritmética. Si se supone nefness en cada fibra
geométrica y curvatura semipositiva, su Corolario~6.4 da una
desigualdad entre volumen e intersección; la positividad de esta
última produce secciones pequeñas \cite{moriwaki-volumen}.
Ser nef en las fibras significa tener grado no negativo en sus
curvas; ser aritméticamente nef exige además grados no negativos
en las curvas horizontales y semipositividad de la curvatura.
Una asintótica para datos fijos necesita además un control uniforme
antes de utilizarse en una familia de pesos, niveles y componentes.

La cantidad de vectores pequeños y el subespacio que generan
son datos distintos. En $\Z^2$, para $N>0$, consideremos
\begin{equation}
 \|(x,y)\|^2=e^{-2N}x^2+e^{2N}y^2.
 \label{eq:geom-secciones-ejemplo}
\end{equation}
Todo vector de norma a lo sumo uno tiene $y=0$, mientras que hay
$2\lfloor e^N\rfloor+1$ posibilidades para $x$.
El determinante de Gram es uno: la pendiente media duplicada
es cero. El conteo crece exponencialmente y los vectores contados
generan una sola recta. A, p.~28, utiliza este ejemplo para
explicar por qué la construcción añade información de generación,
rango y descenso a la información de volumen.

\subsection{Promediar por componentes con sus rangos}

Las componentes de una cubierta contribuyen grados y rangos
diferentes. El promedio relevante para la pendiente está ponderado
por los rangos. El siguiente lema de A, p.~29, conserva esa
distinción.

\begin{lema}[Comparación de promedios]
Sea $\mathcal Z$ finito y no vacío. Supongamos que $h_Z,s_Z\in\R$,
$c_s,c_c,X\ge0$ y $0\le\rho,z<1$ satisfacen
\[
 h_Z+c_s\ge0,\qquad
 (1-\rho)\max(0,s_Z+c_s)\le h_Z+c_s+c_c,\qquad
 \frac1{|\mathcal Z|}\sum_Zs_Z\ge-X.
\]
Sean $d_0,b_T>0$ y números $d_Z,r_Z>0$ tales que
\[
 d_0\le d_Z\le\frac{d_0}{1-\rho},
 \qquad
 b_Td_Z\le r_Z\le\frac{b_Td_Z}{1-z}.
\]
Entonces
\begin{equation}
 \frac{\sum_Zr_Zh_Z}{\sum_Zr_Z}
       \ge-X-(z+2\rho)c_s-c_c.
 \label{eq:geom-promedio}
\end{equation}
\label{lem:geom-promedio}
\end{lema}

\begin{proof}
Pongamos $\theta=(1-z)(1-\rho)$ y
$w_Z=r_Z/\sum_Wr_W$.
Las cotas de rango implican $w_Z\ge\theta/|\mathcal Z|$.
Como $h_Z+c_s\ge0$, resulta
\[
 \sum_Zw_Z(h_Z+c_s)
 \ge\frac{\theta}{|\mathcal Z|}\sum_Z(h_Z+c_s).
\]
La segunda hipótesis y $\max(0,s_Z+c_s)\ge s_Z+c_s$ dan
\[
 \frac1{|\mathcal Z|}\sum_Z(h_Z+c_s)
       \ge(1-\rho)(c_s-X)-c_c.
\]
Al restar $c_s$ se obtiene
\[
 \sum_Zw_Zh_Z
 \ge-\theta(1-\rho)X
       -[1-\theta(1-\rho)]c_s-\theta c_c.
\]
Se tiene $0<\theta\le1$ y
$1-(1-z)(1-\rho)^2\le z+2\rho$.
Estas dos desigualdades prueban la conclusión.
\end{proof}

En la aplicación de A, $b_T=\binom{T-1}{q}$, donde $q$ es una
dimensión geométrica y $T$ una potencia auxiliar. La entrada
de grado por componente tiene la forma
\begin{equation}
 d_Z^{\rm ar}\ge b_Td_Z
 \bigl[Th_Z-3(q+1)|h_Z|-A_0\log(2+|h_Z|)-R_0\bigr].
 \label{eq:geom-grado-componente}
\end{equation}
La construcción aporta las restantes cotas de $T$, $A_0$, $R_0$
y las alturas. El lema explica el coste
$T(z+2\rho)c_s$ al convertir un promedio de alturas en una
estimación ponderada. No autoriza a reemplazar esas alturas
por una cota inferior puntual ausente de las hipótesis.
Tampoco estima por sí solo los términos absoluto y logarítmico
de \eqref{eq:geom-grado-componente}.

\subsection{Mapas de multiplicación y un complemento pequeño}

Sobre una cubierta, A, pp.~28--29, compara espacios de secciones
con pesos cercanos mediante mapas integrales de multiplicación.
La pérdida se puede aislar mediante el siguiente argumento.

\begin{lema}[Transferencia con generación controlada]
Sean $E,F$ retículos hermitianos sobre $\OK$ de rangos positivos.
Sea $I$ un conjunto finito no vacío, y supongamos que
\[
 \Phi:E\longrightarrow F^{\oplus I},\qquad
 \Phi(s)=(\phi_i(s))_{i\in I},
\]
es integral e inyectiva, y que, en cada inmersión $\sigma$,
\begin{equation}
 \sum_{i\in I}\|\phi_i(s)\|_\sigma^2
                         \le e^{2b}\|s\|_\sigma^2
 \quad(b\ge0).
 \label{eq:geom-multiplicacion-norma}
\end{equation}
Supongamos además, para $P>0$ y $A_0\ge0$, que
\[
 \mu(E)\ge-A_0,\qquad
 |I|\operatorname{rk}F\le(1+b/P)\operatorname{rk}E,
\]
y que $F_K$ está generado por un mismo conjunto de vectores de
$F$ cuyas normas son a lo sumo $e^P$ en todas las inmersiones.
Entonces
\begin{equation}
                   \mu(F)\ge-A_0-4b.
 \label{eq:geom-multiplicacion-pendiente}
\end{equation}
\label{lem:geom-multiplicacion}
\end{lema}

\begin{proof}
Escribamos $r=\operatorname{rk}E$ y
$n=|I|\operatorname{rk}F$, y sea $S$ la saturación de la imagen
de $\Phi$. La cota de operador $\|\Phi_\sigma\|\le e^b$ limita
el crecimiento del volumen en dimensión $r$ por $e^{br}$.
La comparación de determinantes, con la normalización duplicada,
da $\mu(S)\ge\mu(E)-2b\ge-A_0-2b$.
La saturación solo aumenta el grado de la imagen.

El cociente $Q=F^{\oplus I}/S$ tiene la métrica cociente.
Los generadores indicados de las sumas directas proyectan a
vectores integrales que generan $Q_K$, de norma a lo sumo $e^P$
en cada inmersión. Si su rango es positivo, escojamos entre ellos
una base sobre $K$. La desigualdad de Hadamard acota cada
determinante de Gram por $e^{2P\operatorname{rk}Q}$.
El subretículo generado por esa base tiene pendiente al menos
$-2P$, y su índice en $Q$ solo mejora la cota. Luego
$\mu(Q)\ge-2P$.

La aditividad del grado y la métrica ortogonal de la suma directa
implican
\[
 n\mu(F)
 \ge-r(A_0+2b)-2P(n-r).
\]
Si $n=r$, el último sumando es cero.
En todos los casos $r\le n$ y $(n-r)/n\le b/P$; por tanto,
$\mu(F)\ge-A_0-2b-2b$.
\end{proof}

La escala $P$ corresponde al producto $P\kappa$ de A; aquí se
renombra para no confundir sus factores con otros parámetros.
El argumento usa una cota para la \emph{suma} de las normas al
cuadrado. Cotas separadas $\|\phi_i\|\le e^b$ solo darían
$\|\Phi\|\le |I|^{1/2}e^b$ y añadirían $\log|I|$ a la
comparación de pendientes.
La existencia de esos mapas, la generación integral y el
descenso desde la cubierta siguen siendo entradas geométricas.
La cuenta anterior conserva también la forma más general
\[
 \mu(F)\ge t-\frac bP(|t|+2P)
 \quad\text{si}\quad t\le\mu(E)-2b.
\]
En efecto, la cota previa a simplificar es
$\mu(F)\ge(r/n)t-2P(1-r/n)$.
Para transportar un término principal de signo libre se debe
conservar esta pérdida de rango; el corolario con $4b$ tiene
la hipótesis adicional $\mu(E)\ge-A_0$, $A_0\ge0$.

\subsection{Proyecciones integrales y rango descartado}

El recorte que selecciona tipos o representaciones conserva
una parte del retículo. Su pérdida depende del rango eliminado,
de sus pendientes y del denominador de la proyección.

\begin{lema}[Coste de una proyección ortogonal racional]
Sean $0\ne L\subseteq M$ retículos hermitianos con la métrica
inducida, y $e:L_K\to L_K$ un idempotente ortogonal en todas
las inmersiones. Supongamos que un entero $g\ge1$ satisface
$ge(L)\subseteq L$. Definamos
\[
 L_g=L\cap\operatorname{im}e,\quad
 L_b=L\cap\ker e,\quad
 r_g=\operatorname{rk}L_g,\quad r_b=\operatorname{rk}L_b.
\]
Si $0\le\delta_b\le1/12$ y $Q_0\ge0$ satisfacen
\[
 r_b\le3\delta_b\operatorname{rk}L,\qquad
 \mu_{\max}(M)\le Q_0,\qquad
 \mu(L)\ge-Q_0,\qquad 2\log g\le Q_0,
\]
entonces $L_g\ne0$ y
\begin{equation}
                \mu(L_g)\ge\mu(L)-12\delta_bQ_0.
 \label{eq:geom-proyeccion}
\end{equation}
\label{lem:geom-proyeccion}
\end{lema}

\begin{proof}
La descomposición sobre $K$ da $r=r_g+r_b=\operatorname{rk}L$.
Como $r_b\le r/4$, se tiene $r_g>0$ y
\[
             \frac{r_b}{r_g}
              \le\frac{3\delta_b}{1-3\delta_b}\le4\delta_b.
\]
La proyección $1-e$ induce un isomorfismo
\[
 L/(L_g\oplus L_b)\simeq(1-e)L/L_b.
\]
La inclusión $g(1-e)L\subseteq L_b$ acota el orden de este
cociente por $g^{[K:\Q]r_b}$.
La ortogonalidad y la fórmula de índice dan
\[
 r\mu(L)\le r_g\mu(L_g)+r_b\mu(L_b)+2r_b\log g.
\]
Si $r_b=0$, el cociente es trivial y la afirmación sigue
directamente. Si $r_b>0$, la saturación de $L_b$ en $M$
tiene pendiente no menor que la de $L_b$, de modo que
$\mu(L_b)\le\mu_{\max}(M)\le Q_0$.
Despejar $\mu(L_g)$ produce una pérdida a lo sumo
\[
 \frac{r_b}{r_g}\bigl(\mu(L_b)-\mu(L)+2\log g\bigr)
                    \le4\delta_b(3Q_0).
\]
\end{proof}

Este es el Lema~7.2 de A, p.~30. El entero $g$ controla un índice
finito, aunque $e$ sea ortogonal en el espacio complejo. La
Nota~\ref{sec:proyecciones} ya muestra ese fenómeno en dimensión
dos. La cota del rango descartado tiene efecto cuantitativo
después de acotar también las pendientes y ese índice.

\subsection{Absorber una pérdida exponencial de rango}

A, pp.~28--30, organiza la construcción mediante un retículo
central $L_{\rm core}$, un retículo de tipo $L\subseteq M$
y el retículo retenido $L_g$, con el proyector y las métricas
del Lema~\ref{lem:geom-proyeccion}.
Sean $\mathcal E\ge b\ge1$, $C_1,C_2\ge0$, y supongamos
\[
 \mu(L_{\rm core})\ge-\mathcal M_j-C_2\mathcal E,\qquad
 \mu(L)\ge\mu(L_{\rm core})-C_1\mathcal E.
\]
Aquí $\mathcal M_j$ es la contribución principal del periodo
en esa etapa. Supongamos además
\[
 \mu_{\max}(M)\le x^{A_1}b,\qquad
 |\mathcal M_j|\le x^{A_2}b,\qquad
 2\log g\le b,\qquad
 r_b\le3\delta_b(x)r,
\]
con $x\ge1$, $A_1,A_2\ge0$ y
$\delta_b(x)=\exp(-x^c)$, $c>0$ fijo.
Aplicamos el lema anterior con
\[
 Q_0=(x^{A_1}+x^{A_2}+C_1+C_2+1)\mathcal E.
\]
Si
\begin{equation}
 12\delta_b(x)(x^{A_1}+x^{A_2}+C_1+C_2+1)\le1,
 \label{eq:geom-absorcion-condicion}
\end{equation}
se cumplen sus hipótesis y resulta
\begin{equation}
          \mu(L_g)\ge-\mathcal M_j-(C_1+C_2+1)\mathcal E.
 \label{eq:geom-core-final}
\end{equation}
Para cualquier exponente fijo $a$,
$\log(x^ae^{-x^c})=a\log x-x^c\to-\infty$.
Por ello \eqref{eq:geom-absorcion-condicion} vale por encima
de un umbral común una vez fijados los exponentes y constantes.

Si la construcción suministra además
$\operatorname{lcv}(M_j^{\rm row})/D\le-\mu(L_g)$,
con su normalización de covolumen y rango $D$, se obtiene
\[
 \frac{\operatorname{lcv}(M_j^{\rm row})}{D}
       \le\mathcal M_j+(C_1+C_2+1)\mathcal E.
\]
Esta es la deducción numérica de A, Proposición~7.3.
La cota inferior del retículo central, la cota superior máxima
y la comparación final de retículos tienen funciones distintas
en el argumento.

\subsection{Evaluaciones, jets y pendientes máximas}

Evaluar una sección en un punto conserva su valor.
Un jet de orden $k$ conserva además sus coeficientes de Taylor
hasta grado total $k$. En coordenadas locales de dimensión $n$,
el espacio de jets de una función escalar en un punto liso es
\[
 \mathcal O_{X,p}/\mathfrak m_p^{k+1};
\]
los monomios $z^\alpha$ con $|\alpha|\le k$ forman una base.
Su dimensión es $\binom{n+k}{n}$, por el conteo de las soluciones
$\alpha_1+\cdots+\alpha_n+\alpha_{n+1}=k$ en enteros no negativos.
Para una sección de un fibrado en rectas se usa también una
trivialización de su fibra. Una estructura integral de jets
requiere el modelo y los módulos correspondientes sobre el
anillo de enteros.

La inyectividad de una evaluación en varios puntos o jets
significa que ninguna sección no nula del espacio considerado
se anula con todos los órdenes prescritos. Es una propiedad
global de ese espacio de secciones.
Las cotas de Cauchy de la Nota~\ref{sec:periodos} estiman las
normas de coeficientes; la inyectividad y la integridad requieren
sus propias justificaciones.

\begin{lema}[Cota máxima mediante una evaluación integral]
Sea $j:M\hookrightarrow J$ una aplicación integral e inyectiva
entre retículos hermitianos de rangos positivos. Con la altura
global utilizada en la Nota~\ref{sec:marcos},
\begin{equation}
                  \mu_{\max}(M)\le\mu_{\max}(J)+h(j).
 \label{eq:geom-evaluacion-max}
\end{equation}
\label{lem:geom-evaluacion-max}
\end{lema}

\begin{proof}
Para cada subretículo saturado $S\subseteq M$ de rango positivo,
la comparación de determinantes de $j|_S$ da
$\mu(S)\le\mu(jS)+h(j|_S)$.
Su altura es a lo sumo $h(j)$ porque la norma restringida no
excede la del operador completo.
La saturación de $jS$ en $J$ aumenta su grado, de donde
$\mu(jS)\le\mu_{\max}(J)$.
Tomar el máximo sobre $S$ prueba la desigualdad.
\end{proof}

Este lema describe qué debe aportar una construcción por
evaluaciones para controlar la pendiente máxima. Una evaluación
compleja sobreyectiva, una cota de su norma y una realización
integral son afirmaciones diferentes. También se distinguirá
la sobreyectividad, necesaria para interpolar datos, de la
inyectividad que se acaba de utilizar.

\subsection{Interpolación con jacobiano y pesos globales}

A, Teorema~7.4, p.~31, utiliza una entrada de extensión $L^2$.
Sea $X$ una variedad compleja de dimensión $n\ge1$, con un mapa
localmente biholomorfo $z:X\to\C^n$ y una función de agotamiento
plurisubarmónica suave. La medida $\dd V_z$ es la medida euclídea
transportada por esas coordenadas. Sean $g=(g_1,\ldots,g_n)$
holomorfo, $\varphi,\varphi_2$ plurisubarmónicas suaves,
$r>0$ y $j\ge2$ entero. Las hipótesis incluyen
\begin{equation}
 \begin{gathered}
 \|g(x)\|^2e^{-\varphi(x)}\le r^2\quad(x\in X),\qquad
 Z=\{g=0\}\text{ finito},\\
 J_g(p)=\det(dg/dz)(p)\ne0\quad(p\in Z).
 \end{gathered}
 \label{eq:geom-ot-hip}
\end{equation}
La plurisubarmonicidad suave expresa la semipositividad de la
forma de Levi; el agotamiento permite aproximar $X$ por
subniveles relativamente compactos.

La entrada afirma que, dados $c_p\in\C$, existe una función
holomorfa $F$ tal que $F(p)=c_pJ_g(p)$ y
\begin{equation}
 \int_X|F|^2e^{-\varphi_2-(n+j)\varphi}\dd V_z
 \le C_{\rm OT}\sum_{p\in Z}|c_p|^2
                           e^{-\varphi_2(p)-j\varphi(p)},
 \qquad
 C_{\rm OT}=\frac{280\pi^nr^{2n}}{(n-1)!}.
 \label{eq:geom-ot}
\end{equation}
La existencia holomorfa se recibe como resultado citado.
La teoría de extensión de Ohsawa--Takegoshi proporciona el
contexto analítico; \cite[§4]{demailly-extension} explica el
método de extensión con estimaciones.
La constante $280$ y las normalizaciones de
\eqref{eq:geom-ot} son las de la entrada específica de A.

El jacobiano se entiende mediante una cuenta local.
Pongamos $\rho=\|g\|^2e^{-\varphi}/r^2$.
Cerca de un cero simple, $w=g(z)$ es un cambio de coordenadas y
$\dd V_w=|J_g|^2\dd V_z$.
En la corona $\varepsilon^2/2\le\rho\le\varepsilon^2$,
el integrando
$|J_g|^2e^{-\varphi_2-(n+j)\varphi}\rho^{-n}$
se transforma en
\[
 r^{2n}|w|^{-2n}e^{-\varphi_2-j\varphi}\dd V_w.
\]
La fórmula polar
\begin{equation}
 \int_{a\le|w|^2\le b}|w|^{-2n}\dd V_w
                  =\frac{\pi^n}{(n-1)!}\log(b/a)
 \quad(0<a<b)
 \label{eq:geom-corona}
\end{equation}
se obtiene integrando
$2\pi^nR^{2n-1}/(n-1)!$ contra $R^{-2n}\dd R$.
Para justificar el límite, se reescala
$w=\varepsilon r e^{\varphi(p)/2}u$.
La condición de la corona pasa a
$1/2\le|u|^2e^{\varphi(p)-\varphi(z(w))}\le1$.
Para $\varepsilon$ pequeño, estos dominios quedan dentro de
$1/4\le|u|^2\le2$ y sus indicadoras convergen, salvo en las
esferas frontera de medida cero, a la de
$1/2\le|u|^2\le1$.
Los pesos convergen uniformemente y el integrando queda dominado
en esa corona fija. La convergencia dominada da la masa límite
\[
 \frac{\pi^n\log2}{(n-1)!}
                 r^{2n}e^{-\varphi_2(p)-j\varphi(p)}.
\]
Este cálculo explica el factor $r^{2n}$ y la desaparición de
$n$ en el peso puntual. Para interpolar valores $F(p)=a_p$,
se debe poner $c_p=a_p/J_g(p)$; por tanto, el lado derecho
contiene $|a_p|^2/|J_g(p)|^2$.
La condición global de \eqref{eq:geom-ot-hip} y la elección
de medida forman parte de la estimación.

\subsection{De la interpolación al determinante de evaluación}

\begin{lema}[Gram de una evaluación sobreyectiva]
Sea $\mathcal H$ un espacio de Hilbert, sea $Y$ un espacio
hermitiano de dimensión $q>0$ y sea
$E:\mathcal H\to Y$ lineal y continuo.
Si para todo $c\in Y$ existe $f\in\mathcal H$ con
$Ef=c$ y $\|f\|^2\le C\|c\|^2$, donde $C>0$, entonces
\begin{equation}
                 EE^*\ge C^{-1}I_Y,\qquad
                 \det(EE^*)\ge C^{-q}.
 \label{eq:geom-eval-gram}
\end{equation}
\label{lem:geom-eval-gram}
\end{lema}

\begin{proof}
La aplicación es sobreyectiva y su núcleo es cerrado.
Cada fibra tiene un único elemento de norma mínima, en
$(\ker E)^\perp$. Esos elementos definen una inversa por
la derecha lineal $S:Y\to\mathcal H$, con $\|S\|\le\sqrt C$.
Para $c\in Y$,
\[
 \|c\|^2=\langle ESc,c\rangle
        =\langle Sc,E^*c\rangle
        \le\sqrt C\|c\|\|E^*c\|.
\]
Luego $\langle EE^*c,c\rangle=\|E^*c\|^2
\ge C^{-1}\|c\|^2$.
Todos los autovalores de $EE^*$ son al menos $C^{-1}$;
su producto prueba la desigualdad de determinantes.
\end{proof}

En \eqref{eq:geom-ot}, tomamos como $\mathcal H$ el espacio
de funciones holomorfas de norma $L^2$ finita para el peso
indicado, y como $Y$ el espacio $\C^Z$ con norma
\[
 \|c\|_Y^2=\sum_{p\in Z}|c_p|^2e^{-\varphi_2(p)-j\varphi(p)}.
\]
La evaluación es $E(F)=(F(p)/J_g(p))_{p\in Z}$.
La desigualdad local de valor medio y los pesos suaves positivos
garantizan su continuidad. Para $Z\ne\varnothing$, el lema
da \eqref{eq:geom-eval-gram} con $C=C_{\rm OT}$, calculando
el adjunto y el determinante con esas métricas.
Es el Lema~7.5 de A, pp.~31--32.

La evaluación de este espacio de Bergman es de orden cero.
Una interpolación de jets de orden superior exige otro enunciado.
Además, restringir $E$ a un subespacio de formas automorfas
puede destruir su sobreyectividad: el interpolante construido
en $\mathcal H$ debe pertenecer al modelo que se quiere utilizar.
La comparación aritmética requiere también identificar
el retículo de evaluación, sus índices y todas sus métricas.

A, p.~32, expresa una comparación posterior entre niveles
mediante las pendientes normalizadas
$\mu_0=\mu(R_0^\vee)-\log\kappa_R$ y
$\mu_\sharp=\mu(R^\vee)-\log\kappa_R$:
\begin{equation}
 \mu_\sharp\le
 \mu_0+3\delta_e(\mu_0+h_0+28\lambda)+28\lambda.
 \label{eq:geom-base-sharp}
\end{equation}
Aquí la construcción suministra la misma normalización de
filas $\kappa_R$ en ambos niveles. Si
$\delta_e\ge0$, $\mu_0+h_0+28\lambda\le Q_0$,
$Q_0\ge1$, $3\delta_eQ_0\le1$ y
$0\le\lambda\le\lambda_0$, la pérdida es a lo sumo
$1+28\lambda_0$.
El número $28$ resume comparaciones anteriores; la sola
desigualdad $\log280<28$ no demuestra
\eqref{eq:geom-base-sharp}.

\subsection{Correlaciones por bloques y las dos trazas}

Sea $(A_\alpha:H_\alpha\to\mathcal H_0)_\alpha$ una familia
finita no vacía de aplicaciones entre
espacios hermitianos complejos de dimensiones
$r_\alpha>0$ y $D>0$.
Los bloques
$K_{\alpha\beta}=A_\alpha^*A_\beta$
comparan columnas completas.
Las entradas espectrales de A, Proposición~7.6, p.~32,
y B, pp.~45--46, tienen la forma
\begin{align}
 \|K_{\alpha\alpha}-I_{H_\alpha}\|_{\rm HS}^2
       &\le e_0r_\alpha,
       \label{eq:geom-bloque-diagonal}\\
 \|K_{\alpha\beta}\|_{\rm HS}^2
       &\le(1+\eta_0)\frac{r_\alpha r_\beta}{D_*}
       \quad(\alpha\ne\beta),
       \label{eq:geom-bloque-cruzado}\\
 D&\le(1+\eta_1)D_*,
       \label{eq:geom-dimension}
\end{align}
donde $D_*>0$ y $e_0,\eta_0,\eta_1\ge0$.
Los escalares positivos de la normalización de los periodos
se incorporan en $A_\alpha$ como una elección de coordenadas
complejas auxiliares. Las alturas aritméticas $h_\alpha$
conservan los mapas y métricas originales de la
Nota~\ref{sec:periodos}.
La norma de Hilbert--Schmidt del bloque es independiente de
las bases ortonormales; una cota de entradas individuales
necesitaría contabilizar sus dimensiones.

\begin{lema}[De bloques a trazas]
Pongamos $a=\sqrt{e_0}$, $R=\sum_\alpha r_\alpha$ y
$Q=\sum_\alpha A_\alpha A_\alpha^*$.
Bajo \eqref{eq:geom-bloque-diagonal}--\eqref{eq:geom-dimension},
\begin{equation}
 |\operatorname{Tr}Q-R|\le aR,\qquad
 \operatorname{Tr}(Q^2)
       \le(1+a)^2R+(1+\eta_0)\frac{R^2}{D_*}.
 \label{eq:geom-trazas}
\end{equation}
Si $e_0\le1/4$, las hipótesis del marco de
\eqref{eq:marcos-error-fuentes}--\eqref{eq:marcos-error-destino}
se cumplen con $\gamma_\alpha=1$, $e=e_0$ y
\begin{equation}
 1+\eta=
 \frac{(D/R)(1+a)^2+(1+\eta_0)(1+\eta_1)}{(1-a)^2}.
 \label{eq:geom-eta-marco}
\end{equation}
\label{lem:geom-bloques-trazas}
\end{lema}

\begin{proof}
Sea $E_\alpha=K_{\alpha\alpha}-I$.
Cauchy--Schwarz para la traza implica
\[
 |\operatorname{Tr}E_\alpha|
 \le\sqrt{r_\alpha}\|E_\alpha\|_{\rm HS}
 \le ar_\alpha.
\]
Además,
\[
 \|K_{\alpha\alpha}\|_{\rm HS}^2
 =r_\alpha+2\operatorname{Tr}E_\alpha+\|E_\alpha\|_{\rm HS}^2
 \le(1+a)^2r_\alpha.
\]
Sumar las trazas da la primera cota.
Por ciclicidad,
\[
 \operatorname{Tr}(Q^2)
    =\sum_{\alpha,\beta}
                 \|A_\alpha^*A_\beta\|_{\rm HS}^2.
\]
Los bloques diagonales aportan a lo sumo $(1+a)^2R$;
los cruzados aportan a lo sumo
$(1+\eta_0)\sum_{\alpha\ne\beta}r_\alpha r_\beta/D_*
\le(1+\eta_0)R^2/D_*$.
Como $\operatorname{Tr}Q\ge(1-a)R>0$, multiplicar esta
cota por $D/(\operatorname{Tr}Q)^2$ y usar
\eqref{eq:geom-dimension} prueba la condición del receptor.
La condición de las fuentes se obtiene sumando
\eqref{eq:geom-bloque-diagonal}. La expresión elegida
para $\eta$ es no negativa.
\end{proof}

La comparación utiliza los pesos de rango del teorema de
pendientes: normalizar $Q$ por $D/R$ corresponde a
$p_\alpha=r_\alpha/R$.
Una elección distinta de pesos debe conservar esa
correspondencia antes de aplicar la Nota~\ref{sec:marcos}.
Las cotas espectrales
\eqref{eq:geom-bloque-diagonal}--\eqref{eq:geom-dimension}
son entradas especializadas: su demostración utiliza
estimaciones locales, expansiones convergentes de núcleos,
selección de representaciones y conteos uniformes.

\subsection{Recortes en las fuentes y en el receptor}

En la aplicación, los espacios fuente completos se recortan
antes de usar su filtración de pendientes. La estimación siguiente
reconstruye la última transición de B, p.~45.

\begin{proposicion}[Marco después de un recorte de fuentes]
Supongamos \eqref{eq:geom-bloque-diagonal}--\eqref{eq:geom-dimension},
y escribamos ahora $d_\alpha$ para los rangos completos y
$R_0=\sum_\alpha d_\alpha$.
Sean $V_\alpha:\C^{r_\alpha}\hookrightarrow H_\alpha$
inclusiones isométricas, $B_\alpha=A_\alpha V_\alpha$ y
$R=\sum_\alpha r_\alpha$.
Se omiten las columnas con $r_\alpha=0$.
Supongamos
\[
 0<\sigma\le\frac1{20},\qquad
 \eta_0,\eta_1,\sqrt{e_0}\le\frac{\sigma}{2},\qquad
 R_0\ge\frac{D_*}{\sigma},\qquad
 R\ge(1-\delta_0)R_0,\quad 0\le\delta_0\le\frac14.
\]
Entonces las aplicaciones $B_\alpha$ satisfacen las dos
condiciones del marco con
\begin{equation}
              e=2e_0,\qquad \eta=5\sigma+3\sqrt{\delta_0}.
 \label{eq:geom-marco-recortado}
\end{equation}
\label{prop:geom-compresion}
\end{proposicion}

\begin{proof}
Pongamos $a=\sqrt{e_0}$,
$E_\alpha=A_\alpha^*A_\alpha-I$ y
$P_\alpha=V_\alpha V_\alpha^*$.
La compresión no aumenta la norma de Hilbert--Schmidt, así que
\[
 \sum_\alpha\|B_\alpha^*B_\alpha-I\|_{\rm HS}^2
 \le e_0R_0\le\frac{e_0}{1-\delta_0}R\le2e_0R.
\]
Sea $Q=\sum_\alpha B_\alpha B_\alpha^*$ y $t=\operatorname{Tr}Q$.
Cauchy--Schwarz para la suma de los bloques da
\[
 \left|\sum_\alpha\operatorname{Tr}(P_\alpha E_\alpha)\right|
 \le\sqrt R\sqrt{\sum_\alpha\|E_\alpha\|_{\rm HS}^2}
 \le aR_0.
\]
Por ello $t\ge(1-\delta_0-a)R_0>0$.
Si $Q_0=\sum_\alpha A_\alpha A_\alpha^*$, entonces
$0\le Q\le Q_0$ y
$\operatorname{Tr}(Q^2)\le\operatorname{Tr}(Q_0^2)$.
En efecto, escribir $Q_0=Q+D_0$ con $D_0\ge0$ produce
$2\operatorname{Tr}(QD_0)+\operatorname{Tr}(D_0^2)\ge0$;
no se necesita que los dos operadores conmuten.
El lema anterior implica
\begin{equation}
 \frac{D\operatorname{Tr}(Q^2)}{t^2}
 \le\frac{(1+\eta_1)[1+\eta_0+\sigma(1+a)^2]}
              {(1-\delta_0-a)^2}
 \le\frac{N(\sigma)}{(1-\delta_0-\sigma/2)^2},
 \label{eq:geom-razon-compresion}
\end{equation}
donde
\[
 N(s)=1+2s+\frac74s^2+\frac34s^3+\frac18s^4.
\]
Para $0\le s\le1/20$, $N(s)\le9/8$ y
\[
 (1+4s)(1-s/2)^2-N(s)
 =s\left(1-\frac{11}{2}s+\frac14s^2-\frac18s^3\right)\ge0.
\]
La última desigualdad se sigue, por ejemplo, de
$1-\frac{11}{2}s-\frac18s^3
\ge1-\frac{11}{40}-\frac1{64000}>0$.
Para $0\le d\le1/4$,
\[
 \frac{\dd}{\dd d}\frac{N(s)}{(1-d-s/2)^2}
 \le\frac{2(9/8)}{(29/40)^3}
 =\frac{144000}{24389}<6.
\]
Integrar desde $d=0$ da
\[
 \frac{N(s)}{(1-d-s/2)^2}
 \le1+4s+6d\le1+5s+3\sqrt d.
\]
La última comparación usa $6d\le3\sqrt d$ para $d\le1/4$.
Aplicarla a \eqref{eq:geom-razon-compresion} prueba la
segunda condición del marco.
Finalmente, $2e_0\le\sigma^2/2\le1/800<1/4$, como exige
el teorema de la Nota~\ref{sec:marcos}.
\end{proof}

B, p.~45, suministra para su recorte ramificado la identidad
\[
 r_\alpha=d_\alpha
       \left(1-\frac1{T_\alpha-1}\right)^f.
\]
Esta identidad de dimensiones es una entrada de la representación
local. A partir de ella, si $T_\alpha\ge b\ge3$, la desigualdad
de Bernoulli da
\begin{equation}
 1-\frac{r_\alpha}{d_\alpha}
 \le\frac{f}{T_\alpha-1}\le\frac{2f}{b},
 \qquad
 \frac{R_0-R}{R_0}\le\frac{2f}{b}.
 \label{eq:geom-rango-ramificado}
\end{equation}

El recorte en el receptor tiene otro parámetro.
El espacio espectral completo, de dimensión $D$, se proyecta
mediante una misma aplicación sobre el receptor aritmético
$C\otimes_\sigma\C$, de dimensión $d$.
La construcción debe realizar los mapas $b_\alpha$ mediante
esa proyección y dar
\begin{equation}
                      d\ge(1-\delta_{\rm row})D,
                 \qquad 0\le\delta_{\rm row}<1.
 \label{eq:geom-rango-receptor}
\end{equation}
Este $\delta_{\rm row}$ es el $\delta$ de la Nota~\ref{sec:marcos}.
A, p.~33, recibe además
$\delta_{\rm row}\le\exp(-x^{c_{\rm row}})$, con
$c_{\rm row}>0$ fijo. Por encima de un umbral,
$\delta_{\rm row}\le1/2$ y
$1/(1-\delta_{\rm row})\le2$.
El parámetro $\delta_0$ de \eqref{eq:geom-marco-recortado}
mide fuentes descartadas; $\delta_b$ mide el recorte del
retículo del Lema~\ref{lem:geom-proyeccion}.
Mantenerlos separados permite aplicar cada cota con su rango
y su denominador correctos.

\subsection{Dispersión de pendientes y precisión del marco}

Con rangos fuente retenidos $r_\alpha$ y alturas $h_\alpha$,
pongamos
\begin{align}
 s&=\sum_\alpha\frac{r_\alpha}{R}
                    \bigl(\mu(N_\alpha)-h_\alpha\bigr),
                    \label{eq:geom-media-ajustada}\\
 u&=\max_\alpha\bigl(\mu_{\max}(N_\alpha)-h_\alpha\bigr)-s.
                    \label{eq:geom-dispersion}
\end{align}
Se tiene $u\ge0$: el máximo domina la media ponderada de
$\mu(N_\alpha)-h_\alpha$.
Para $C=M_W^\vee$, la cota de profundidad requerida es
\[
 s-\mu_{\min}(C)=s+\mu_{\max}(M_W)\le H_0.
\]
B, p.~46, emplea una cota $H_0=bx^{C_H}$ con exponente fijo.
La pendiente máxima del retículo de filas y la pendiente media
que se estima mediante el covolumen son entradas distintas.

Hay una cuenta útil para localizar el efecto de los índices.
Si $M'_\alpha\subseteq M_\alpha$ tiene índice finito y
$N_\alpha=\mathfrak a_\alpha^{-1}M'_\alpha$, pongamos
\[
 \iota_\alpha=
 \frac{2}{[K:\Q]r_\alpha}\log\#(M_\alpha/M'_\alpha).
\]
El giro por el ideal desplaza todas las pendientes por un
mismo número $g_\alpha$. La Nota~\ref{sec:periodos} da
\[
 \mu(N_\alpha)=\mu(M_\alpha)+g_\alpha-\iota_\alpha,
 \qquad
 \mu_{\max}(N_\alpha)\le\mu_{\max}(M_\alpha)+g_\alpha.
\]
Por consiguiente,
\begin{equation}
 \mu_{\max}(N_\alpha)-\mu(N_\alpha)
 \le\mu_{\max}(M_\alpha)-\mu(M_\alpha)+\iota_\alpha.
 \label{eq:geom-dispersion-indice}
\end{equation}
Si $a_\alpha=\mu(N_\alpha)-h_\alpha$, se deduce también
\[
 u\le
 \max_\alpha\{\mu_{\max}(M_\alpha)-\mu(M_\alpha)+\iota_\alpha\}
       +\max_\alpha a_\alpha-\min_\alpha a_\alpha.
\]
Esta desigualdad identifica los términos que necesitan cotas;
no afirma que cada uno sea pequeño sin utilizar la construcción.

La pérdida de la transferencia es
\begin{equation}
 \mathcal E_{\rm marco}
 =\frac{4u\log(1+H_0/u)+(\eta+4\sqrt e)H_0}
        {1-\delta_{\rm row}},
 \label{eq:geom-presupuesto-marco}
\end{equation}
con el primer término igual a cero para $u=0$.
Para $H_0\ge0$, la identidad
\[
 u\log(1+H_0/u)=\int_0^{H_0}\frac{u}{u+t}\dd t
 \quad(u>0)
\]
muestra que la expresión es creciente en $u$ y se prolonga
continuamente por cero en $u=0$.
Una cota $u\le U$ puede por tanto sustituirse en ese término.

Los otros errores se multiplican por $H_0$.
Por ejemplo, para lograr
$\sigma H_0\le bx^{-A}$ cuando $H_0\le bx^{C_H}$,
es suficiente escoger
\begin{equation}
 \sigma=x^{-(C_H+A)},\qquad
 q_{\rm col}\ge C_0\Gamma^2/\sigma,
 \qquad
 d_\alpha\ge\frac{D_*}{C_0\Gamma^2}.
 \label{eq:geom-cardinalidad}
\end{equation}
La segunda y tercera condiciones dan
$R_0\ge D_*/\sigma$.
Si se escoge un $\sigma$ menor, la cardinalidad debe ajustarse
a ese nuevo valor.
Para $C_H+A>0$, aumentar el umbral asegura $\sigma\le1/20$.

A, p.~32, escribe $\Gamma=\exp(C_\Gamma m^5fL)$.
Las cotas espectrales proporcionadas en A y B son
$\eta_0,\eta_1,\sqrt{e_0}\le\exp(-x^{\varepsilon/8})$.
Con $\varepsilon>0$ fijo, esa cantidad es menor que toda
potencia inversa fija de $x$ por encima de un umbral.
La existencia de suficientes columnas admisibles con las
dimensiones de \eqref{eq:geom-cardinalidad} es una entrada
adicional de selección y conteo.

La pérdida $\sqrt{\delta_0}$ requiere también atención.
Si $\delta_0\le2f/b$, entonces
\[
 \frac{\sqrt{\delta_0}H_0}{b}
       \le x^{C_H}\sqrt{\frac{2f}{b}}.
\]
En particular, con $f\le x^{C_f}$ y $b=\exp(x^\varepsilon)$,
el lado derecho es a lo sumo
$\sqrt2\,x^{C_H+C_f/2}\exp(-x^\varepsilon/2)$
y tiende a cero más deprisa que cualquier potencia inversa fija.
La afirmación está normalizada por $b$; no dice que el
producto sin normalizar tienda a cero.

Primero se fijan los exponentes de las cotas de alturas y
pendientes. Después se escoge la precisión que requieren
sus productos con los errores del marco y, finalmente,
la cantidad de columnas y los umbrales espectrales.
Un máximo de finitísimos umbrales ya uniformes permite
ensamblar las estimaciones. Umbrales individuales para cada
columna de una familia que crece con $x$ no dan ese resultado.

\subsection{Entradas geométricas y aritméticas delimitadas}

Las pruebas anteriores separan los mecanismos finitos de la
construcción de sus datos. La tabla precisa los puntos en que
el estudio todavía recibe resultados especializados.

\begin{center}
\small
\begin{tabularx}{\textwidth}{@{}>{\raggedright\arraybackslash}p{0.22\textwidth}Y>{\raggedright\arraybackslash}p{0.18\textwidth}@{}}
\toprule
Objetivo & Datos que deben construirse & Localizador\\
\midrule
Pendiente media inferior &
Modelos y métricas de frontera; alturas y grados de componentes;
mapas de multiplicación integrales; generadores pequeños y descenso. &
A, pp.~28--30; B, pp.~33--40.\\
Pendiente máxima &
Evaluaciones o comparaciones que controlen el mismo retículo,
con inyectividad, integridad y alturas uniformes. &
A, pp.~28--32; B, pp.~39--40.\\
Determinante de evaluación &
Extensión holomorfa con los pesos globales; paso al espacio
automorfo utilizado y a su retículo integral. &
A, pp.~31--32.\\
Dos trazas del marco &
Cotas diagonales y cruzadas, convergencia de expansiones,
comparación de dimensiones y suficientes columnas admisibles. &
A, pp.~32--33; B, pp.~44--46.\\
Costes finitos &
Denominadores de proyectores, longitudes locales finitas
y sumabilidad de la contribución cohomológica. &
A, pp.~33--41; B, §5.\\
\bottomrule
\end{tabularx}
\end{center}

En particular, la métrica de Hodge natural puede tener
singularidades en la frontera. B, pp.~37--38, describe una
corrección por un potencial $U/2$, regularizaciones y una
comparación de autointersecciones. Esos pasos justifican
el uso de resultados formulados para métricas suaves.
El control de los límites y del promedio de $U$ pertenece
a la construcción; no se deduce del solo nombre del
teorema de Hilbert--Samuel.

En los lugares finitos, la entrada de jets utilizada en A, §6,
tiene la hipótesis $p\ge5$. Los primos $2$ y $3$ reciben
comparaciones integrales distintas, descritas en A, pp.~33--34.
La suma de defectos locales de A, (7.9), exige longitudes
cohomológicas finitas y una serie no negativa convergente.
Esa definición de defecto no incorpora el factor dos:
este aparece al convertir el coste a grado y pendiente.

Con estas entradas, la geometría y las dos trazas suministran
las hipótesis que la Nota~\ref{sec:periodos} inserta en su
balance. Producir las ventanas aritméticas exige además
controlar los denominadores y escoger simultáneamente los
parámetros de selección. Las desigualdades de esta nota
muestran qué cantidades y qué uniformidad debe conservar
esa elección.

\subsection{Concordancia localizada con Lean}
\label{subsec:geometria-lean}

El cotejo utiliza la revisión
\href{https://github.com/CocoSgt/riemann-hypothesis-lean/tree/7cbf7a0c32fca1badf05931879a37c3f8a2b03b4}{\texttt{7cbf7a0}}
del mismo espejo \cite{lean}.
Se leyeron 19 archivos completos y pasajes delimitados de
otros cuatro. El registro conserva las 23 rutas, hashes,
intervalos y distinción entre interfaces y pruebas.
Un archivo reutilizado de la Nota~\ref{sec:periodos} se
volvió a leer y a identificar por sus bytes.

\begingroup
\raggedright
\begin{enumerate}
\item \textbf{Componentes y descenso de cubiertas.}
\href{https://github.com/CocoSgt/riemann-hypothesis-lean/blob/7cbf7a0c32fca1badf05931879a37c3f8a2b03b4/src/Development/AbelianVarietyTheory/Definitions/Divisors/CovolXOfCovol.lean}{\texttt{CovolXOfCovol}}
define los predicados de rangos, mapas, generadores, grados
y alturas. Su cota de generadores es simultánea en todas
las inmersiones, con escala \texttt{PP * kap}.
La prueba
\href{https://github.com/CocoSgt/riemann-hypothesis-lean/blob/7cbf7a0c32fca1badf05931879a37c3f8a2b03b4/src/Development/AbelianVarietyTheory/Proofs/DivisorProofs/HtMapLeOfPoints.lean}{\texttt{descent}}
reúne los mapas en una suma ortogonal y obtiene
$\mu(F)\ge t-\epsilon(|t|+2\Lambda)$ bajo sus hipótesis
de rango y norma. El Lema~\ref{lem:geom-multiplicacion}
desarrolla aquí una variante con cota inicial
$\mu(E)\ge-A_0$, $A_0\ge0$, y pérdida $4b$,
a partir de la desigualdad previa de grados.
Se leyeron también
\href{https://github.com/CocoSgt/riemann-hypothesis-lean/blob/7cbf7a0c32fca1badf05931879a37c3f8a2b03b4/src/Formalism/MeanHeightDescent/AverageDegreeBounds/SumComponents.lean}{\texttt{sumComponents'}}
y
\href{https://github.com/CocoSgt/riemann-hypothesis-lean/blob/7cbf7a0c32fca1badf05931879a37c3f8a2b03b4/src/Development/AbelianVarietyTheory/Proofs/DivisorProofs/UniformLowerBoundOfSums.lean}{\texttt{uLow\_of}},
que conservan los rangos ponderados y las entradas de
altura. El Lema~\ref{lem:geom-promedio} aísla una parte
de ese argumento; no se presenta como transcripción literal
de toda la prueba formal de componentes.

\item \textbf{Recorte y absorción.}
En
\href{https://github.com/CocoSgt/riemann-hypothesis-lean/blob/7cbf7a0c32fca1badf05931879a37c3f8a2b03b4/src/Development/AbelianVarietyTheory/Proofs/DivisorProofs/MajAtAddRight.lean}{\texttt{MajAtAddRight}},
\texttt{slope\_good\_ge} conserva la pérdida $12\delta Q$,
el denominador integral y la ortogonalidad en todas las
inmersiones. En
\href{https://github.com/CocoSgt/riemann-hypothesis-lean/blob/7cbf7a0c32fca1badf05931879a37c3f8a2b03b4/src/Development/AbelianVarietyTheory/Proofs/DivisorProofs/CovolumeFromGramAndBounds.lean}{\texttt{CovolumeFromGramAndBounds}},
\texttt{lcv\_le} realiza la absorción y
\texttt{covolX\_of\_steps} reúne umbrales uniformes.
El punto de entrada
\href{https://github.com/CocoSgt/riemann-hypothesis-lean/blob/7cbf7a0c32fca1badf05931879a37c3f8a2b03b4/src/Development/AutomorphicArithmetic/ThetaCorrespondence/ThetaPairLifts/HodgeBundleWeights/TopCovolSlim2Use.lean}{\texttt{topCovolSlim2\_use\_alt}}
recibe explícitamente resultados geométricos, de volumen,
de comparación integral y de separación de puntos.
Se cotejó su ensamblaje; no se reconstruyó la clausura de
esas entradas.

\item \textbf{Interpolación.}
\href{https://github.com/CocoSgt/riemann-hypothesis-lean/blob/7cbf7a0c32fca1badf05931879a37c3f8a2b03b4/src/ReciprocityTheory/L2DbarEstimates/JacobianInterpolationStatement.lean}{\texttt{JacobianInterpolationStatement}}
prescribe valores $c(p)J_g(p)$ con el peso de
\eqref{eq:geom-ot}; su entero $j$ es un exponente del peso.
La prueba interna
\href{https://github.com/CocoSgt/riemann-hypothesis-lean/blob/7cbf7a0c32fca1badf05931879a37c3f8a2b03b4/src/ReciprocityTheory/L2DbarEstimates/LeviForms/Ot0.lean}{\texttt{I.ot0}}
recibe cuatro entradas de solución, coronas, datos locales y corte.
El teorema exterior \texttt{ot0} de
\href{https://github.com/CocoSgt/riemann-hypothesis-lean/blob/7cbf7a0c32fca1badf05931879a37c3f8a2b03b4/src/ReciprocityTheory/L2DbarEstimates/NormFibrations/Hormander.lean}{\texttt{NormFibrations/Hormander}}
construye esas entradas mediante resultados importados;
su firma no las exige como hipótesis adicionales.
Se leyeron ambas pruebas, sin auditar todos esos imports.
La deducción abstracta de Gram se demuestra aquí como
argumento matemático; no se identificó en este alcance
un único teorema Lean que fuera su transcripción exacta.

\item \textbf{Defectos colectivos.}
\href{https://github.com/CocoSgt/riemann-hypothesis-lean/blob/7cbf7a0c32fca1badf05931879a37c3f8a2b03b4/src/Formalism/SizeWindowBrackets/TraceSquareAndDiagBounds.lean}{\texttt{TraceSquareAndDiagBounds}}
define el operador, los defectos y la compresión.
Las pruebas están en
\href{https://github.com/CocoSgt/riemann-hypothesis-lean/blob/7cbf7a0c32fca1badf05931879a37c3f8a2b03b4/src/Formalism/SizeWindowBrackets/ChainSlopeWindows/TraceConjTransposeMulSelf.lean}{\texttt{TraceConjTransposeMulSelf}}:
\texttt{trace\_Qm\_sq}, \texttt{compressDiagDefectBound}
y \texttt{compressTraceSquareBound}, entre otras.
Se mantienen $e=2e_0$ y $\eta=5\sigma+3\sqrt{\delta_0}$.
La prueba de la Proposición~\ref{prop:geom-compresion}
desarrolla una deducción suficiente de esas constantes,
sin afirmar identidad paso a paso con su prueba Lean.

\item \textbf{Los mismos mapas y sus dos rangos.}
Las estructuras de
\href{https://github.com/CocoSgt/riemann-hypothesis-lean/blob/7cbf7a0c32fca1badf05931879a37c3f8a2b03b4/src/Formalism/LocalToGlobal/GlobalTheory/FamRec.lean}{\texttt{FamRec}}
y las pruebas seleccionadas de
\href{https://github.com/CocoSgt/riemann-hypothesis-lean/blob/7cbf7a0c32fca1badf05931879a37c3f8a2b03b4/src/Formalism/LocalToGlobal/GlobalTheory/FamRecP.lean}{\texttt{FamRecP}}
identifican las matrices comprimidas, la proyección al
receptor y los mapas aritméticos.
\texttt{del} mide el déficit de filas y \texttt{deF}
el de las fuentes.
\href{https://github.com/CocoSgt/riemann-hypothesis-lean/blob/7cbf7a0c32fca1badf05931879a37c3f8a2b03b4/src/Formalism/LocalToGlobal/SiegelLatticeTheory/FlagStabilizerTensors/LinkSFinal.lean}{\texttt{LinkSFinal}}
ensambla el enlace espectral desde estimaciones locales,
de núcleos y de conteo.
\texttt{LinkS\_at} sigue siendo una implicación bajo la
comparabilidad de $d_\alpha\Gamma^2/D_*$;
\texttt{LinkS} fija la cota
$\exp(-x^{\varepsilon/8})$ para $\eta_0$, $\eta_1$
y $\sqrt{e_0}$.
Los campos de una estructura de datos no constituyen
por sí solos una demostración de que esos datos existen.
\end{enumerate}
\endgroup

Los hashes recalculados coinciden con los archivos recuperados
en la revisión fijada. No se ejecutó ni compiló Lean, no se
auditó el cierre de axiomas y dependencias y no se verificó
íntegramente la identidad del espejo con la distribución
original de Alpöge. La finitud y sumabilidad cohomológicas,
el paso completo de Bergman al retículo automorfo y la elección
final de parámetros quedan fuera del alcance demostrado aquí.

\subsection{Ejercicios con soluciones}

\begin{ejercicio}\label{ej:geom-pequenos}
En $\Z^2$, con $N>0$, considere la norma al cuadrado
$e^{-2N}x^2+e^{4N}y^2$.
Determine los vectores de norma a lo sumo uno, su subespacio
generado, el covolumen y la pendiente media duplicada.
\end{ejercicio}

\paragraph{Solución del Ejercicio~\ref{ej:geom-pequenos}.}
Como $e^{4N}>1$, se debe tener $y=0$.
La condición restante es $|x|\le e^N$.
Hay $2\lfloor e^N\rfloor+1$ vectores y generan la recta
$\Q(1,0)$.
Las longitudes ortogonales de la base son $e^{-N}$ y
$e^{2N}$; por tanto, el covolumen es $e^N$ y
$\mu=-2\log(e^N)/2=-N$.
El número de vectores pequeños crece mientras la pendiente
media tiende a $-\infty$.

\begin{ejercicio}\label{ej:geom-jets}
En $M=\Z[z]_{\le3}$ compare la aplicación de coeficientes
de Taylor en cero con
$f\mapsto(f(0),f'(0),f''(0),f'''(0))$.
Calcule sus índices en $\Z^4$ y el efecto de reducir módulo
dos. En $\Z[z]_{\le1}$, calcule además el índice de
$f\mapsto(f(0),f(2))$.
\end{ejercicio}

\paragraph{Solución del Ejercicio~\ref{ej:geom-jets}.}
Si $f=a_0+a_1z+a_2z^2+a_3z^3$, el jet normalizado es
\[
 \bigl(f(0),f'(0),f''(0)/2!,f'''(0)/3!\bigr)
                         =(a_0,a_1,a_2,a_3).
\]
Es un isomorfismo integral y su índice es uno.
Las derivadas ordinarias tienen matriz
$\operatorname{diag}(1,1,2,6)$, con índice $12$.
Módulo dos desaparecen sus dos últimas coordenadas;
los coeficientes normalizados conservan las cuatro.
Su integridad en este ejemplo procede del retículo de
polinomios enteros.

Para $f=a+bz$, la evaluación da $(a,a+2b)$.
Su imagen es
$\{(u,v)\in\Z^2:v-u\in2\Z\}$, de índice dos.
La matriz tiene determinante dos y es un isomorfismo
sobre $\Q$ o $\C$.
La interpolación compleja no elimina esa condición de
congruencia sobre los valores enteros.

\begin{ejercicio}\label{ej:geom-recorte}
Tome $q$ copias de $A_\alpha=I_{\C^D}$.
Restrinja todas a un mismo subespacio de dimensión
$(1-\delta_0)D$, entero positivo.
Calcule el defecto de isometría de cada restricción y
el menor $\eta$ que satisface la desigualdad de la
segunda traza en dimensión $D$.
\end{ejercicio}

\paragraph{Solución del Ejercicio~\ref{ej:geom-recorte}.}
Cada restricción es isométrica, así que el defecto diagonal
es cero. Si $P$ es la proyección sobre el subespacio retenido,
el operador colectivo es $Q=qP$ y
\[
 \operatorname{Tr}Q=q(1-\delta_0)D,\qquad
 \operatorname{Tr}(Q^2)=q^2(1-\delta_0)D.
\]
Por ello
$D\operatorname{Tr}(Q^2)/(\operatorname{Tr}Q)^2
=1/(1-\delta_0)$ y el mínimo es
$\eta=\delta_0/(1-\delta_0)$.
Aumentar $q$ conserva ese cociente: las direcciones
eliminadas siguen siendo las mismas.

\begin{ejercicio}\label{ej:geom-uniforme}
Suponga $H_0\le bx^{C_H}$,
$d_\alpha\ge D_*/(C_0\Gamma^2)$,
$C_H\ge0$ y $A>0$.
Elija $\sigma$ y una cardinalidad suficiente para obtener
$\sigma H_0\le bx^{-A}$ y $R_0\ge D_*/\sigma$.
Explique también por qué la convergencia a cero de
$f_k(x)=\mathbf1_{\{x<k\}}$ para cada $k$ fijo no da
convergencia uniforme en $k\in\N$.
\end{ejercicio}

\paragraph{Solución del Ejercicio~\ref{ej:geom-uniforme}.}
Se puede escoger
$\sigma=x^{-(C_H+A)}$ y un número entero de columnas
\[
 q_{\rm col}\ge C_0\Gamma^2x^{C_H+A}.
\]
Entonces $\sigma H_0\le bx^{-A}$ y
$R_0\ge q_{\rm col}D_*/(C_0\Gamma^2)\ge D_*/\sigma$.
Las cotas espectrales y $\sigma\le1/20$ se imponen con
un umbral común posterior.
Un $\sigma$ menor requiere volver a exigir
$q_{\rm col}\ge C_0\Gamma^2/\sigma$.

Para cada $k$ fijo, $f_k(x)=0$ cuando $x\ge k$.
Sin embargo, dado cualquier $x$, se puede elegir un entero
$k>x$, y entonces $f_k(x)=1$.
Por tanto, $\sup_k f_k(x)=1$ para todo $x$.
Este ejemplo distingue los cuantificadores
$\forall k\,\exists x_k$ y $\exists x_0\,\forall k$ que
intervienen en una estimación uniforme.




